\begin{claim}{D.2.24}{Proved}{A primary operation excludes the remaining dangerous unit pattern}
Put $B=K(\FF_2,2)$, $X=\Sigma B$, $A=QX$, $Y=\Sigma^2B$,
and $T=\Omega^2Y$. Retain the labels and classes of
\textup{D.2.17--D.2.22}, in particular
\[
 k=\{a,b\},\qquad p_3=f+ac+a^3,\qquad
 d_6=b^2,\quad e_6=b^2+k,\quad f_6=p_3.
 \tag{D.2.24a}
\]
Let $d\in H_5(B;\FF_2)$ be dual to
$w=\operatorname{Sq}^2\operatorname{Sq}^1u$ and zero on
$u\operatorname{Sq}^1u$, and write $D=\sigma d\in H_6X$.
For every ordinary pointed map $g:X\to T$, the coefficient
of $e_6$ in $g_*D$ is zero. Equivalently, its coefficient
on $k$ in the primitive basis $b^2,k,p_3$ is zero.

Consequently the pattern $G_*D=b^2+k$ constructed through
dimension nine in \textup{D.2.17--D.2.20} cannot extend over
the unit ten-skeleton. This conclusion applies to every map
with that pattern, independently of its earlier cell fillings.

For an arbitrary ordinary map $G:A\to T$, there is
$s\in\FF_2$ such that, on the unit classes and the indicated
extra product,
\[
 G_*D=s p_3,\qquad G_*E=0,\qquad G_*(A_0^2)=0,
 \quad E=\sigma e,\quad A_0=\sigma a.
 \tag{D.2.24b}
\]
Its once-suspended adjunct $\Sigma A\to\Omega Y$ therefore
has zero coefficient on the old commutator
$\bar a\bar b+\bar b\bar a$, where $\bar a=\sigma a$ and
$\bar b=\sigma b$ are the loop-homology letters, in its
entire degree-seven homology image. This is a statement about
ordinary maps and the first attaching-map test; it does not
assert persistence through later nullification stages.
\end{claim}
\begin{proof}
All homology and cohomology in this proof has coefficients
$\FF_2$. Let $F=\fib(T\to B)$ be the three-connected fiber
used above. Its homology injects into that of $T$.
Write $e_6^*$ and $J_1^*$ for the dual basis vectors to
$e_6$ and $J_1=\{a,d\}$ in $H_6F$ and $H_8F$.
The complete eighteen-dimensional table
\textup{(D.2.22a6)} gives the exact cohomology identity
\[
 \Psi_F:=\operatorname{Sq}^4(e_6^*)+
                   \operatorname{Sq}^2(J_1^*)=0
                   \quad\hbox{in }H^{10}F.
 \tag{D.2.24c}
\]
Indeed the only degree-ten basis vectors whose fourth square
has $e_6$ coefficient one are $B_4,B_5$. These are exactly
the two basis vectors whose second square has $J_1$
coefficient one. Every other basis vector contributes zero
to both functionals. Thus this is an identity in the full
cohomology group, with no untested complementary summand.
The table follows from the opened CLM formulas
\cite[Chapter III, Theorems 1.2--1.4]{I_CohenLadaMay1976};
no homotopy-group presentation or secondary attachment
coefficient is needed for \textup{(D.2.24c)}.

Write $v=\operatorname{Sq}^1u$, $w=\operatorname{Sq}^2v$,
and $t=\operatorname{Sq}^4w$ in the polynomial presentation
of $H^*B$. Serre's generator theorem and Cartan give
\[
\begin{aligned}
 \operatorname{Sq}^4w&=t,&
 \operatorname{Sq}^4(uv)&=u^2w+v^3,\\
 \operatorname{Sq}^2(u^2v)&=u^2w+v^3,&
 \operatorname{Sq}^2(uw)&=u^2w+v^3.
\end{aligned}
 \tag{D.2.24d}
\]
The degree-five cohomology basis is $w,uv$, and the
degree-seven basis is $u^2v,uw$. The complete degree-nine
basis is $u^3v,u^2w,v^3,t$; see
\cite[\S2, nos.8--9]{I_Serre1953}.
Let $t_9\in H_9B$ be dual to $t$ and zero on the other
three basis monomials. By suspension compatibility,
the class $T_{10}=\sigma t_9\in H_{10}X$ satisfies
\[
 \operatorname{Sq}_*^4T_{10}=D,
 \qquad \operatorname{Sq}_*^2T_{10}=0.
 \tag{D.2.24e}
\]

Every pointed $g:X\to T$ lifts to $F$: its composite
with $T\to K(\FF_2,2)$ is null because $H^2X=0$.
Use the same letter for this lift. Naturality and the
ordinary homology--cohomology pairing now give
\[
\begin{split}
 0
 &=\langle\Psi_F,g_*T_{10}\rangle\\
 &=\langle e_6^*,g_*\operatorname{Sq}_*^4T_{10}\rangle
   +\langle J_1^*,g_*\operatorname{Sq}_*^2T_{10}\rangle\\
 &=\langle e_6^*,g_*D\rangle.
\end{split}
 \tag{D.2.24f}
\]
This proves the universal assertion. It uses neither an
$H$-map hypothesis nor naturality of Dyer--Lashof operations
under an arbitrary ordinary map. The same argument already
applies on a unit ten-skeleton: choose a lift of $T_{10}$
to its homology; inclusion is an isomorphism in the lower
degrees occurring in \textup{(D.2.24e)}. Since
$b^2+k=e_6$, the previously constructed dangerous pattern
cannot extend over this skeleton.

For completeness, we recall the low-degree constraints needed
for the last assertion. For $G:A\to T$, \textup{D.2.16}
and the primitive basis in degree five give $G_*(\sigma c)=0$.
Also $G_*A_0=0$ and $G_*(\sigma b)=s a^2$ for some $s$.
The complete primitive bases in degrees six and seven are
\[
 \langle b^2,k,p_3\rangle,
 \qquad
 \langle\xi b,\{a,c\},P,Q\rangle,
 \tag{D.2.24g}
\]
where $P,Q$ are as in \textup{D.2.22}. Naturality of
$\beta$, $\operatorname{Sq}_*^2$, and
$\operatorname{Sq}_*^3$ on the unit classes gives,
on writing $G_*D=d_1b^2+d_2k+d_3p_3$ and
$G_*E=e_1b^2+e_2k+e_3p_3$,
\[
 d_1=e_1,\qquad d_2=e_2=\epsilon,\qquad
 d_1+d_3=s,\qquad e_1=e_3.
 \tag{D.2.24h}
\]
The exact second-Bockstein calculation of
\textup{(D.2.17e)} has
$\beta_2[b^2]=\beta_2[k]=[\xi a]$ and $[p_3]=0$.
On the source $\beta_2[D]=[\sigma c]$.
Since $G_*(\sigma c)=0$, naturality gives
$d_1=d_2=\epsilon$. Hence
\[
 G_*D=\epsilon(b^2+k)+(s+\epsilon)p_3,
 \qquad
 G_*E=\epsilon(b^2+k+p_3).
 \tag{D.2.24i}
\]
Apply \textup{(D.2.24f)} to the unit restriction of $G$.
It forces $\epsilon=0$.

There is no additional image on $A_0^2$. The source class
$R_0=Q^4A_0$ satisfies $\beta R_0=A_0^2$ and
$\operatorname{Sq}_*^2R_0=0$. The second identity excludes
a $\xi b$ coefficient in $G_*R_0$, so its Bockstein is
a multiple of $p_3$. Applying the second square, and using
$\operatorname{Sq}_*^2p_3=a^2$, excludes that multiple.
Thus $G_*(A_0^2)=0$, completing \textup{(D.2.24b)}.

Finally $H_6A$ has basis $D,E,A_0^2$. Homology suspension
from $T$ to $\Omega Y$ sends $k$ to the old tensor
commutator, sends $b^2$ to zero, and sends $p_3$ to
the linear $f$ label. Formula \textup{(D.2.24b)} therefore
has no commutator coefficient on any degree-seven input
of the once-suspended adjunct. This proves the stated
scope without making an assertion about subsequent cones.
\end{proof}
